
\def\PUDcodigo{ECA225}

\def\PUDcodacad{04507.42}

\def\PUDdisciplina{Dispositivos Periféricos}

\def\PUDsemestre{S8}

\def\PUDhoras{80}

%NOVOS ---------------------------------------------------
\def\PUDhorasTeoria{40}
\def\PUDhorasPratica{40}

\def\PUDinterECA{ 
\hyperref[PUD:ECA219]{Instrumentação (04507.25)}

\hyperref[PUD:ECA213]{Microcontroladores (04507.28)}

\hyperref[PUD:ECA233]{Controle I (04507.31)}

\hyperref[PUD:ECA221]{CLP (04507.35)} 
}

\def\PUDatuaisECA{
- Gravação em banco de dados \\
- Internet das coisas - IoT (Esp8266/NodeMCU) 
}

\def\PUDrecursos{Projetor multimídia; Quadro branco e pincel; Aulas em laboratório.}

%----------------------------------------------------------

\def\PUDcreditos{4}

\def\PUDprerequisito{ \hyperref[PUD:ECA213]{ECA213} }

\def\PUDpreacad{ \hyperref[PUD:ECA213]{Microcontroladores (04507.28)} }

\def\PUDobjetivos{Conhecer os princípais dispositivos periféricos aplicados à Robótica e à Sistemas Embarcados utilizando Microcontroladores.  Projetar sistemas e/ou aplicações com tais dispositivos.}

\def\PUDementa{Introdução aos dispositivos de entrada e saída aplicados à Microcontroladores; 
Tipos de atuadores e sensores aplicados à Robótica;  Tipos de dispositivos aplicados à Sistemas Embarcados;  Projetos de sistemas e/ou aplicações utilizando sistemas embarcados.}

\def\PUDconteudo{ & Introdução aos dispositivos de entrada e saída aplicados à Microcontroladores. 
\\ 
 & Tipos de dispositivos atuadores e sensores aplicados à Robótica (Motor de Passo, Servo Motor, Motor CC, driver de potência, Encoder, transmissor e receptor infravermelho, transmissor e receptor por ultrasom, bumper, sensor fim de curso, entre outros). 
\\ 
 & Tipos de dispositivos aplicados à Sistemas Embarcados (LDR, LM35, LCD, RTC, LED RGB, LED de alto brilho, teclado matricial, memória EEPROM interna e externa, entre outros). 
\\ 
 & Comunicação entre dispositivos (Bluetooth, Wifi, Zigbee, rádio frequência, entre outras). Gravação em banco de dados e IoT (Esp8266/NodeMCU).
\\ 
 & Desenvolvimento de projetos utilizando dispositivos periféricos. }

\def\PUDmetodo{Aulas expositórias que podem ser teóricas e/ou práticas, onde as práticas no laboratório serão marcadas no decorrer da disciplina.}

\def\PUDavalia{A avaliação será feita com aplicação de provas teóricas e/ou práticas, além da possibilidade de inclusão de trabalhos e seminários no decorrer da disciplina.}

\def\PUDbibbasica{ &  \href{www.biblioteca.ifce.edu.br/index.asp?codigo_sophia=23871}{Souza}, José David de.  Desbravando o PIC.  12º ed.  São Paulo, SP: Érica, 2008.  268p.  ISBN: 9788571948679.
\\ 
 & \href{www.biblioteca.ifce.edu.br/index.asp?codigo_sophia=24245}{Thomazini}, D; Urbano, P.  Sensores industriais: fundamentos e aplicações: funcionamento e especificações, tipos de sensores e aplicações na indústria.  4º ed. São Paulo, SP: Érica, 2007.  222p.  ISBN: 9788536500713.
\\ 
 %&  MADRID, Marconi Kolm.  Curso sobre robôs industriais.  Fortaleza (CE): Universidade Federal do Ceará - UFC, 1992.  92 p. }
 & \href{www.biblioteca.ifce.edu.br/index.asp?codigo_sophia=24331}{Nicolosi}, Denis E. C.  Laboratório de Microcontroladores: Família 8051, Treino de instruções, hardware e software.  5º ed.  São Paulo, SP: Érica, 2006.  206p.  ISBN: 8571948712. 
 }

\def\PUDbibcomplementar{ & \href{www.biblioteca.ifce.edu.br/index.asp?codigo_sophia=39837}{Lathi}, B.  P.  Sinais e Sistemas Lineares.  2º ed.  Porto Alegre, RS: Bookman, 2007.  856p.  ISBN: 9788560031139.
\\ 
 &  \href{www.biblioteca.ifce.edu.br/index.asp?codigo_sophia=24070}{Costa}, Cesar da.  Projetando controladores digitais com FPGA: aprenda a projetar o seu próprio Controlador Lógico Programável (CLP) utilizando FPGA de baixo custo.  São Paulo, SP: Novatec, 2006.  159p.  ISBN: 8575220888.
\\ 
 &  \href{www.biblioteca.ifce.edu.br/index.asp?codigo_sophia=24581}{Vasconcelos}, Laércio.  Hardware na prática: construindo e configurando micros de 32 e 64 bits single core, dual core e quad core para usuários, técnicos e estudantes.  2º ed.  Rio de Janeiro, RJ: Laércio Vasconcelos Computação, 2007.  724p.  ISBN: 9788586770074. \\
 & \href{www.biblioteca.ifce.edu.br/index.asp?codigo_sophia=57234}{Aguirre}, Luis Antonio.  Fundamentos de Instrumentação.  E-book.  Pearson.  354p.  ISBN: 9788581431833.
\\
 & \href{www.biblioteca.ifce.edu.br/index.asp?codigo_sophia=24365}{Arrabaça}, Devair Aparecido.  Eletrônica de potência: conversores de energia (CA/CC): teoria, prática e simulação.  São Paulo, SP: Érica, 2011.  334p.  ISBN: 9788536503714. }

\def\PUDautor{Pedro Pedrosa Rebouças Filho}

\def\PUDdata{2013-06-06}

\def\PUDrevisao{2}

\def\PUDrevisor{Fábio Timbó Brito}

\def\PUDdatarevisao{2019-05-15}

\PUDcorpo
